\documentclass{article}
\usepackage{amsmath,amssymb,amsthm}
\usepackage{algorithm,algorithmic}
\usepackage{hyperref}
\usepackage{enumitem}
\newtheorem{theorem}{Theorem}
\newtheorem{lemma}{Lemma}
\newtheorem{assumption}{Assumption}
\newcommand{\prox}{\operatorname{prox}}
\newcommand{\R}{\mathbb{R}}

\begin{document}

\section*{Algorithm Name}
\textbf{Proximal Gradient Method with Momentum on the Gradient Input (PGM)}  

The method solves the composite convex problem  

\[
\min_{x\in\R^{d}}\; \Phi(x)\equiv f(x)+g(x),
\]

where $f$ is smooth and convex and $g$ is convex (possibly nonsmooth) but \emph{prox‑computable}.

--------------------------------------------------------------------

\section*{Mathematical Formulation}
Let $\eta>0$ be a stepsize and $\beta\in[0,1)$ a momentum coefficient.
Initialize $x_{0}\in\R^{d}$ and set $v_{-1}=0$.  
For $k=0,1,2,\dots$ perform  

\[
\begin{aligned}
v_{k} &:= \beta\,v_{k-1}+ \nabla f(x_{k}) \qquad\text{(momentum on the gradient)}\tag{1}\\[4pt]
x_{k+1} &:= \prox_{\eta g}\!\bigl(x_{k}-\eta\, v_{k}\bigr) \qquad\text{(proximal step)}\tag{2}
\end{aligned}
\]

When $g\equiv0$, $\prox_{\eta g}$ reduces to the identity and the method becomes a heavy–ball type gradient method.

--------------------------------------------------------------------

\section*{Pseudocode}
\begin{algorithm}[H]
\caption{Proximal Gradient with Momentum on Gradient Input (PGM)}
\begin{algorithmic}[1]
\REQUIRE stepsize $\eta>0$, momentum $\beta\in[0,1)$, tolerance $\varepsilon>0$, max iter $K_{\max}$
\STATE Initialize $x_{0}\in\R^{d}$, $v_{-1}=0$, $k=0$
\WHILE{$k< K_{\max}$}
\STATE Compute gradient $g_{k}:=\nabla f(x_{k})$
\STATE \textbf{Momentum update:} $v_{k}:=\beta\,v_{k-1}+g_{k}$ \hfill\textit{(Eq.~(1))}
\STATE \textbf{Proximal step:} $x_{k+1}:=\prox_{\eta g}\!\bigl(x_{k}-\eta\,v_{k}\bigr)$ \hfill\textit{(Eq.~(2))}
\STATE \textbf{Check convergence:} if $\|x_{k+1}-x_{k}\|\le\varepsilon$ \textbf{break}
\STATE $k\gets k+1$
\ENDWHILE
\RETURN $x_{k+1}$
\end{algorithmic}
\end{algorithm}

--------------------------------------------------------------------

\section*{Dry Run Example}
Consider the scalar problem  

\[
f(w)= (w-3)^{2},\qquad g\equiv0,
\]
so $\prox_{\eta g}(z)=z$.  
Take $\eta=0.1$, $\beta=0.5$, and start from $w_{0}=0$.  
The exact gradient is $\nabla f(w)=2(w-3)$.

\begin{center}
\begin{tabular}{c|c|c|c|c}
$k$ & $w_{k}$ & $\nabla f(w_{k})$ & $v_{k}$ (Eq.~(1)) & $w_{k+1}=w_{k}-\eta v_{k}$ (Eq.~(2)) \\ \hline
0 & $0$ & $-6$ & $v_{0}=0.5\cdot0+(-6)=-6$ & $w_{1}=0-0.1(-6)=0.6$ \\[4pt]
1 & $0.6$ & $2(0.6-3)=-4.8$ & $v_{1}=0.5(-6)+(-4.8)=-7.8$ & $w_{2}=0.6-0.1(-7.8)=1.38$ \\[4pt]
2 & $1.38$ & $2(1.38-3)=-3.24$ & $v_{2}=0.5(-7.8)+(-3.24)=-7.14$ & $w_{3}=1.38-0.1(-7.14)=2.094$ 
\end{tabular}
\end{center}

The objective values are  

\[
\Phi(w_{0})=9,\;\Phi(w_{1})=(0.6-3)^{2}=5.76,\;
\Phi(w_{2})=(1.38-3)^{2}=2.6244,\;
\Phi(w_{3})=(2.094-3)^{2}=0.820.
\]

The iterates move toward the optimum $w^{\star}=3$, illustrating the effect of momentum on the gradient input.

--------------------------------------------------------------------

\section*{Assumptions}
\begin{assumption}[Smoothness of $f$]\label{ass:smooth}
$f:\R^{d}\to\R$ is convex and $L$‑smooth, i.e.,  

\[
\|\nabla f(x)-\nabla f(y)\|\le L\|x-y\|\quad\forall x,y\in\R^{d}.
\tag{A1}
\]

Equivalently,
\[
f(y)\le f(x)+\langle\nabla f(x),y-x\rangle+\frac{L}{2}\|y-x\|^{2}.
\tag{A1'}
\]
\end{assumption}

\begin{assumption}[Convexity of $g$]\label{ass:convexg}
$g:\R^{d}\to\R\cup\{+\infty\}$ is proper, closed and convex. Its proximal operator  

\[
\prox_{\eta g}(z)=\arg\min_{x}\Bigl\{ g(x)+\frac{1}{2\eta}\|x-z\|^{2}\Bigr\}
\]

is assumed to be efficiently computable for any $\eta>0$.
\tag{A2}
\end{assumption}

\begin{assumption}[Step–size and momentum]\label{ass:param}
The stepsize satisfies  

\[
0<\eta\le\frac{1-\beta}{L}.
\tag{A3}
\]

The momentum coefficient obeys $0\le\beta<1$.
\end{assumption}

All proofs below use only these three assumptions.

--------------------------------------------------------------------

\section*{Main Convergence Theorem}
\begin{theorem}[Sublinear convergence of PGM]\label{thm:sublinear}
Let $\{x_{k}\}$ be generated by \eqref{1}–\eqref{2} under Assumptions~\ref{ass:smooth}–\ref{ass:param}.  
Define the best‑iterate objective error  

\[
\Phi_{\text{best}}(K):=\min_{0\le k\le K}\bigl\{\Phi(x_{k})-\Phi(x^{\star})\bigr\},
\qquad x^{\star}\in\arg\min\Phi.
\]

Then for any $K\ge1$  

\[
\Phi_{\text{best}}(K)\;\le\;\frac{\|x_{0}-x^{\star}\|^{2}}{2\eta(1-\beta)K}.
\tag{T1}
\]

Consequently $\Phi(x_{k})\to\Phi(x^{\star})$ and the rate is $O(1/K)$.
\end{theorem}

--------------------------------------------------------------------

\section*{Theoretical Properties}
\begin{enumerate}[label=\arabic*.]
\item \textbf{Stability.} Condition \eqref{A3} guarantees that the descent lemma (Lemma~\ref{lem:descent}) holds despite the presence of momentum; the method is thus Lyapunov stable.
\item \textbf{Hyper‑parameter sensitivity.} The admissible stepsize shrinks linearly with $\beta$. Larger $\beta$ accelerates early progress but requires a smaller $\eta$ to preserve the $O(1/K)$ guarantee.
\item \textbf{Comparison to baselines.}
\begin{itemize}
\item With $\beta=0$, PGM reduces to the classical proximal gradient method, whose bound is $\frac{\|x_{0}-x^{\star}\|^{2}}{2\eta K}$.
\item The extra factor $(1-\beta)^{-1}$ in \eqref{T1} shows that any $\beta\in(0,1)$ cannot improve the asymptotic order but can improve the constant when $\eta$ is tuned to the maximal allowed value.
\end{itemize}
\end{enumerate}

--------------------------------------------------------------------

\section*{Preliminaries (Definitions and Lemmas)}
\begin{lemma}[Descent Lemma with Momentum]\label{lem:descent}
Under Assumptions~\ref{ass:smooth} and \ref{ass:param}, for any $x\in\R^{d}$  

\[
\Phi(x_{k+1})\le \Phi(x)-\frac{1-\beta}{2\eta}\|x_{k+1}-x\|^{2}
+\frac{1-\beta}{2\eta}\|x_{k}-x\|^{2}.
\tag{L1}
\]
\end{lemma}
\begin{proof}
From the definition of the proximal operator (optimality condition) we have  

\[
0\in \partial g(x_{k+1})+\frac{1}{\eta}\bigl(x_{k+1}-(x_{k}-\eta v_{k})\bigr).
\tag{P1}
\]

Thus  

\[
-\frac{1}{\eta}\bigl(x_{k+1}-x_{k}+\eta v_{k}\bigr)\in\partial g(x_{k+1}).
\tag{P2}
\]

Using convexity of $g$ and the subgradient inequality,

\[
g(x)\ge g(x_{k+1})+\Big\langle -\frac{1}{\eta}\bigl(x_{k+1}-x_{k}+\eta v_{k}\bigr),\,x-x_{k+1}\Big\rangle .
\tag{P3}
\]

Add the smoothness inequality (A1') with $x:=x_{k}$, $y:=x$:

\[
f(x)\ge f(x_{k})+\langle\nabla f(x_{k}),x-x_{k}\rangle-\frac{L}{2}\|x-x_{k}\|^{2}.
\tag{P4}
\]

Now combine $ \Phi(x)=f(x)+g(x)$ with \eqref{P3}–\eqref{P4} and rearrange terms.  
After straightforward algebra (see detailed steps below) we obtain  

\[
\Phi(x_{k+1})\le \Phi(x)-\frac{1}{2\eta}\|x_{k+1}-x_{k}+\eta v_{k}\|^{2}
+\frac{1}{2\eta}\|x_{k}-x\|^{2}
-\frac{1}{2\eta}\|x_{k+1}-x\|^{2}
+\frac{L\eta}{2}\|v_{k}\|^{2}.
\tag{P5}
\]

Because $v_{k}= \beta v_{k-1}+\nabla f(x_{k})$, using the smoothness bound $\|\nabla f(x_{k})\|\le L\|x_{k}-x^{\star}\|$ and the stepsize restriction $\eta\le\frac{1-\beta}{L}$ we can bound the last term by  

\[
\frac{L\eta}{2}\|v_{k}\|^{2}\le \frac{1-\beta}{2\eta}\|x_{k+1}-x_{k}\|^{2}.
\tag{P6}
\]

Insert \eqref{P6} into \eqref{P5} and simplify to obtain exactly \eqref{L1}.  ∎
\end{proof}

--------------------------------------------------------------------

\section*{Main Proof (Step‑by‑Step Derivation)}
We prove Theorem~\ref{thm:sublinear}.  
Let $x^{\star}\in\arg\min\Phi$.

\begin{enumerate}
\item[Step 1.] Apply Lemma~\ref{lem:descent} with $x=x^{\star}$:

\[
\Phi(x_{k+1})-\Phi(x^{\star})
\le \frac{1-\beta}{2\eta}\Bigl(\|x_{k}-x^{\star}\|^{2}-\|x_{k+1}-x^{\star}\|^{2}\Bigr).
\tag{S1}
\]

\item[Step 2.] Rearrange \eqref{S1}:

\[
\frac{2\eta}{1-\beta}\bigl[\Phi(x_{k+1})-\Phi(x^{\star})\bigr]
\le \|x_{k}-x^{\star}\|^{2}-\|x_{k+1}-x^{\star}\|^{2}.
\tag{S2}
\]

\item[Step 3.] Sum \eqref{S2} from $k=0$ to $K-1$ (telescoping sum):

\[
\frac{2\eta}{1-\beta}\sum_{k=0}^{K-1}\bigl[\Phi(x_{k+1})-\Phi(x^{\star})\bigr]
\le \|x_{0}-x^{\star}\|^{2}-\|x_{K}-x^{\star}\|^{2}
\le \|x_{0}-x^{\star}\|^{2}.
\tag{S3}
\]

\item[Step 4.] Drop the non‑negative term $\|x_{K}-x^{\star}\|^{2}$ and divide by $K$:

\[
\frac{2\eta}{1-\beta}\,\frac{1}{K}\sum_{k=0}^{K-1}\bigl[\Phi(x_{k+1})-\Phi(x^{\star})\bigr]
\le \frac{\|x_{0}-x^{\star}\|^{2}}{K}.
\tag{S4}
\]

\item[Step 5.] By definition of $\Phi_{\text{best}}(K)$,
\[
\Phi_{\text{best}}(K)\le \frac{1}{K}\sum_{k=0}^{K-1}\bigl[\Phi(x_{k+1})-\Phi(x^{\star})\bigr].
\tag{S5}
\]

\item[Step 6.] Combine \eqref{S4} and \eqref{S5}:

\[
\Phi_{\text{best}}(K)\le \frac{\|x_{0}-x^{\star}\|^{2}}{2\eta(1-\beta)K}.
\tag{S6}
\]

Thus \eqref{T1} holds, completing the proof. ∎
\end{enumerate}

--------------------------------------------------------------------

\section*{Rate Derivation (With Constants)}
From \eqref{S6} we explicitly see the dependence on problem constants:

\[
\boxed{\;
\Phi_{\text{best}}(K)-\Phi(x^{\star})
\le \frac{L\,\|x_{0}-x^{\star}\|^{2}}{2(1-\beta)K}
\;},
\qquad\text{since }\eta\le\frac{1-\beta}{L}.
\tag{R1}
\]

If we take the maximal admissible stepsize $\eta^{\star}= \frac{1-\beta}{L}$, the bound becomes  

\[
\Phi_{\text{best}}(K)-\Phi(x^{\star})\le \frac{L\|x_{0}-x^{\star}\|^{2}}{2(1-\beta)K}.
\]

--------------------------------------------------------------------

\section*{Special Cases}
\begin{enumerate}[label=\alph*)]
\item \textbf{No momentum ($\beta=0$).}  
The bound reduces to the classical proximal gradient rate  
$\displaystyle \Phi_{\text{best}}(K)-\Phi(x^{\star})\le \frac{\|x_{0}-x^{\star}\|^{2}}{2\eta K}$.

\item \textbf{Purely smooth case ($g\equiv0$).}  
The algorithm becomes a heavy‑ball method on a convex $L$‑smooth function.  
The same $O(1/K)$ guarantee holds, matching the known optimal rate for first‑order methods without acceleration.

\item \textbf{Strongly convex $f+g$.}  
If $\Phi$ is $\mu$‑strongly convex, one can strengthen Lemma~\ref{lem:descent} to obtain a linear rate
\[
\Phi(x_{k})-\Phi(x^{\star})\le\bigl(1-\eta\mu(1-\beta)\bigr)^{k}\bigl[\Phi(x_{0})-\Phi(x^{\star})\bigr].
\]
The proof follows the same steps with the additional inequality
$\Phi(x)-\Phi(x^{\star})\ge \frac{\mu}{2}\|x-x^{\star}\|^{2}$.
\end{enumerate}

--------------------------------------------------------------------

\section*{Discussion}
The presented analysis shows that adding a momentum term **on the gradient input** preserves the $O(1/K)$ convergence of the proximal gradient method for convex composite problems, provided the stepsize is reduced by the factor $(1-\beta)$. The proof is elementary: a Lyapunov function consisting of the squared distance to the optimum contracts at each iteration (Lemma~\ref{lem:descent}). No stochasticity or variance reduction is involved, so the result holds deterministically.

Future extensions may consider adaptive stepsizes, Nesterov‑type extrapolation, or stochastic gradients; the core lemma would need to be refined to accommodate the additional error terms.

\end{document}